\clearpage\fancyhead[L]{\small JUNO.}\SourceLeaf{397}
\begin{center}\textbf{GAUSS AN BODE.}\par\vspace{4pt}\hrule\vspace{4pt}\fontsize{10}{12}\selectfont Astronomisches Jahrbuch für 1819. Seite 219..221. Berlin 1816.\par\vspace{4pt}\hrule\vspace{4pt}\end{center}
\begin{flushright}Göttingen den 10. August 1816.\end{flushright}

\par\indent Ihrem Wunsche zufolge, habe ich das Vergnügen, Ihnen, verehrtester Freund, beigehend Ephemeriden für die nächste Erscheinung der Pallas und Juno zu übersenden, die ich mit einigen Anmerkungen begleite.

\par\indent Die Ephemeride für die Pallas ist von Hrn \textsc{Westphal} aus Schwerin berechnet, welcher sich hier gegenwärtig mit Eifer der Astronomie widmet. Es ist dabei — wie immer seit 5 Jahren — auf die Störungen nach meiner Theorie Rücksicht genommen, und daher die genaueste Uebereinstimmung zu erwarten. Die nächste Opposition hat Hr. \textsc{Westphal}, wie folgt, im Voraus berechnet:
\begin{center}\fontsize{10}{12}\selectfont
1817 Juli 11. $8^{\mathrm u}\;45^{\mathrm m}\;5^{\mathrm s}$ Mittlere Zeit in Göttingen\par
\setlength{\tabcolsep}{3pt}\begin{tabular}{lcr}
Wahre Länge & $.\;.\;.\;.\;.\;.\;.\;.\;.\;.\;.\;.\;.\;.$ & $289^\circ\;3'\;22''$\\
Geocentrische Breite & $.\;.\;.\;.\;.\;.\;.\;.\;.\;.\;.$ & $43\;33\;1$\\
\end{tabular}\end{center}

\par\indent An Licht wird die Pallas im nächsten Jahre sehr schwach sein, und nur 10—11. Grösse haben. Folgendes ist die Uebersicht der Lichtstärke in den letzten und der künftigen Opposition:
\begin{center}\fontsize{10}{12}\selectfont\setlength{\tabcolsep}{11pt}\begin{tabular}{lr|lr}
1812 & $0.01666$ & 1816 & $0.05997$\\
1813 & $0.01475$ & 1817 & $0.01289$\\
1814 & $0.05476$ & 1818 & $0.02122$\\
\end{tabular}\end{center}

\Needspace{9\baselineskip}\par\indent Die Juno habe ich selbst zweimal am Kreismikrometer beobachtet:
\begin{center}\fontsize{10}{12}\selectfont\setlength{\tabcolsep}{5pt}\begin{tabular}{l|c|c|c}
1816 & Mittlere Zeit & Gerade Aufsteig. & Südl. Abweich.\\\hline
Juni 12 & $10^{\mathrm u}\;47^{\mathrm m}\;27^{\mathrm s}5$ & $251^\circ\;46'\;8''0$ & $3^\circ\;42'\;4''7$\\
13 & $10\;27\;34.7$ & $251\;33\;47.6$ & $3\;41\;11.7$\\
\end{tabular}\end{center}

\Needspace{11\baselineskip}\par\indent Bei dem so ungewöhnlich schlechten Wetter, welches im vorigen Frühjahr herrschte, konnten nicht mehr Beobachtungen gemacht werden. Ich habe daraus folgende Resultate für die Opposition abgeleitet:
\begin{center}\fontsize{10}{12}\selectfont
1816 Juni 3. $4^{\mathrm u}\;20^{\mathrm m}\;20^{\mathrm s}$ Mittlere Zeit in Göttingen\par
\setlength{\tabcolsep}{3pt}\begin{tabular}{lcr}
Wahre Länge & $.\;.\;.\;.\;.\;.\;.\;.\;.\;.\;.\;.$ & $252^\circ\;51'\;39''2$\\
Geocentrische Breite & $.\;.\;.\;.\;.\;.\;.\;.\;.$ & $18\;34\;3.1$ Nördl.\\
\end{tabular}\end{center}

\Needspace{9\baselineskip}\par\indent Hr. \textsc{Enke}, welcher nachher diese Rechnung wiederholt und dabei auch die Wilnaer Beobachtungen mit benutzt hat, fand
\begin{center}\fontsize{10}{12}\selectfont
Juni 3 \quad $4^{\mathrm u}\;17^{\mathrm m}\;31^{\mathrm s}$ \quad Mittlere Zeit in Göttingen\par
$252^\circ\;51'\;32''6$\par
$18\;34\;14.6$
\end{center}

\Needspace{15\baselineskip}\par\indent Aus der Combination der vier letzten Oppositionen leitete ich folgende Elemente ab:
\SourceLeaf{398}\begin{center}\fontsize{9}{11}\selectfont\setlength{\tabcolsep}{2pt}\renewcommand{\arraystretch}{1.15}\begin{tabular}{lcr}
Mittlere Länge 1817 Sept. 6. $0^{\mathrm u}$ in Göttingen & $.\;.$ & $8^\circ\;51'\;41''$\\
Sonnennähe & $.\;.\;.\;.\;.\;.\;.\;.\;.\;.\;.\;.\;.\;.$ & $53\;16\;46$\\
Knoten & $.\;.\;.\;.\;.\;.\;.\;.\;.\;.\;.\;.\;.\;.\;.\;.$ & $171\;9\;41$\\
Neigung & $.\;.\;.\;.\;.\;.\;.\;.\;.\;.\;.\;.\;.\;.\;.$ & $13\;3\;49$\\
Excentricitätswinkel & $.\;.\;.\;.\;.\;.\;.\;.\;.\;.\;.$ & $14\;46\;0$\\
Tägliche mittlere tropische Bewegung & $.\;.\;.\;.$ & $812''40187$\\
Logarithm der halben grossen Axe & $.\;.\;.\;.\;.\;.$ & $0.4269223$\\
\end{tabular}\end{center}

\par\indent (Es ist wohl überflüssig zu bemerken, dass diese Elemente blos zu dem Zweck berechnet sind, zu der Berechnung der nächsten Ephemeride zu dienen: so lange die Störungen noch nicht entwickelt sind, ist dies Verfahren das bequemste und genaueste.)
\par\indent Nach diesen Elementen hat Hr. Prof. \textsc{Harding} eine Ephemeride für die nächste Erscheinung berechnet, von welcher ich Ihnen hier eine Abschrift beilege. Dieselbe Rechnung ist auch von Hrn. \textsc{Posselt} ausgeführt, dessen Ephemeride von der gegenwärtigen überall nur in den Secunden abweichend in Hrn. \textsc{von Lindenau}'s Zeitschrift abgedruckt wird. Die nächste Opposition fällt nach Hrn. Professor \textsc{Harding}'s Rechnung aus obigen Elementen.
\begin{center}\fontsize{10}{12}\selectfont
1817 Sept. 6. $4^{\mathrm u}\;10^{\mathrm m}\;30^{\mathrm s}$ mittlere Zeit in Göttingen.\par
\setlength{\tabcolsep}{3pt}\begin{tabular}{lcr}
Wahre Länge & $.\;.\;.\;.\;.\;.\;.\;.\;.\;.\;.\;.\;.$ & $343^\circ\;37'\;36''$\\
Geocentrische Breite & $.\;.\;.\;.\;.\;.\;.\;.\;.\;.\;.$ & $3\;6\;9$ Nördl.\\
Die Lichtstärke 1816 & $.\;.\;.\;.\;.\;.\;.\;.\;.\;.\;.$ & $0.01747$\\
\hfill1817 & $.\;.\;.\;.\;.\;.\;.\;.\;.\;.\;.$ & $0.11398$\\
\end{tabular}\end{center}

\par\indent Erlauben Sie mir noch ein Paar Bemerkungen über veränderliche Sterne. Dem von Herrn Prof. \textsc{Harding} gefundenen in der Jungfrau legt Hr. \textsc{Koch} im Jahrbuche 1818 eine Periode von 295 Tagen bei, ohne anzugeben, worauf sich dieselbe gründet: allein mit den Beobachtungen des Hrn. Prof. \textsc{Harding} verträgt sich diese Periode nicht. — Einen zweiten veränderlichen Stern hat Hr. Prof. \textsc{Harding} im Jahre 1811 im Wassermann südlich unter $\omega^2$ bemerkt, der von der sechsten Grösse allmählich, und zwar in ungefähr $2\frac12$ Monaten zur völligen Unsichtbarkeit übergeht, dessen Lichtwechsel überhaupt aber sich in keine regelmässige Periode zu fügen scheint. Seine mittlere Stellung 1800 war
\begin{center}\fontsize{10}{12}\selectfont\setlength{\tabcolsep}{3pt}\begin{tabular}{lcr}
Gerade Aufsteigung & $.\;.\;.\;.\;.\;.\;.\;.\;.\;.\;.$ & $353^\circ\;21'\;48''5$\\
Jährliche. Ver. & $.\;.\;.\;.\;.\;.\;.\;.\;.\;.\;.\;.\;.\;.\;.$ & $+46.56$\\
Südliche Abweichung & $.\;.\;.\;.\;.\;.\;.\;.\;.\;.$ & $16^\circ\;23'\;13''4$\\
Jährliche Abnahme & $.\;.\;.\;.\;.\;.\;.\;.\;.\;.\;.\;.\;.\;.$ & $19.88$.\\
\end{tabular}\end{center}
